\documentclass[12pt,a4paper]{article}
\usepackage[T2A]{fontenc}
\usepackage[utf8]{inputenc}
\usepackage[russian]{babel}
\usepackage{amsmath,amssymb,amsfonts}
\usepackage{geometry}
\usepackage{titlesec}
\usepackage{enumitem}
\usepackage{booktabs}
\usepackage{xcolor}
\usepackage{hyperref}
\usepackage{mathtools}
\usepackage{tikz}
\usepackage{pgfplots}
\pgfplotsset{compat=1.18}
\usetikzlibrary{calc, arrows.meta}
\usetikzlibrary{patterns, fillbetween, intersections}
\usepackage{fancyhdr}
\geometry{left=2.5cm,right=2.5cm,top=2.5cm,bottom=2.5cm}
\pagestyle{fancy}
\fancyhf{}
\fancyhead[L]{Машинное обучение 1 --- Контрольная 2024 --- Лашкевич Злата БЭАД245}
\fancyfoot[R]{\thepage}
\renewcommand{\headrulewidth}{0.4pt}
\renewcommand{\footrulewidth}{0pt}
\setlength{\headheight}{14.5pt}
\newcommand{\defeq}{\overset{\mathrm{def}}{=\joinrel=}}

\begin{document}
\selectlanguage{russian}

\section*{Задача 1.(2.5 балла)}

Представьте, что вам предложили обучать линейный классификатор на следующий
функционал (назовём его \emph{сферическим}):
\[
Q(w) = \frac{1}{\ell}\sum_{i=1}^{\ell}
\log\bigl(1 + \exp(-\beta\, y_i \cos\theta(w, x_i))\bigr),
\]
где $\theta(w, x_i)$ "--- угол между векторами $w$ и $x_i$,
$\beta > 0$ "--- скалярный гиперпараметр.

Ответьте на вопросы о свойствах данного функционала (в каждом пункте
ожидается, что вы дадите обоснование своему ответу):

\begin{enumerate}
    \item Легко заметить, что здесь функция потерь очень похожа на
    логистическую. Попробуйте максимально свести сферический функционал
    к логистическому "--- так, чтобы в нём тоже появилось скалярное
    произведение весов на признаки. Проанализируйте полученную формулу:
    в чём ключевое отличие сферического функционала от логистического?

    \item Проведём более глубокое исследование. Известно, что при обучении
    градиентным спуском линейной модели на логистическую функцию потерь
    есть проблема. Рассмотрим линейно разделимую выборку с бинарной
    целевой переменной. На каком векторе весов будет достигаться минимум
    функционала с логистической функцией потерь? Решается ли эта
    проблема в сферическом функционале?

    \item Как параметр $\beta$ влияет на поведение сферического
    функционала? Как будет вести себя функционал при очень больших и при
    очень маленьких значениях $\beta$?
\end{enumerate}

\subsection*{1.}

Вспомним формулу для косинуса через скалярное
произведение:
\[
\cos\theta(w, x_i) = \frac{\langle w, x_i\rangle}{\|w\|\cdot\|x_i\|}
\]
Подставим в функционал:
\[
Q(w) = \frac{1}{\ell}\sum_{i=1}^{\ell}
\log\!\left(1 + \exp\!\left(
-\beta\, y_i \cdot \frac{\langle w, x_i\rangle}{\|w\|\cdot\|x_i\|}
\right)\right)
\]
А стандартный логистический функционал:
\[
Q_{\mathrm{log}}(w) = \frac{1}{\ell}\sum_{i=1}^{\ell}
\log\bigl(1 + \exp(-y_i\langle w, x_i\rangle)\bigr)
\]
Получается, что в сферическом функционале вместо скалярного произведения
$\langle w, x_i\rangle$ стоит его \emph{нормированная версия}
$\dfrac{\langle w, x_i\rangle}{\|w\|\cdot\|x_i\|}$, домноженная на
константу $\beta$\\[1mm]
Таким образом, сферический функционал можно записать как логистическую функцию
потерь, применённую к нормализованным векторам:
\[
Q(w) = \frac{1}{\ell}\sum_{i=1}^{\ell}
\log\!\left(1 + \exp\!\left(
-y_i \cdot \beta \cdot \left\langle
\frac{w}{\|w\|},\; \frac{x_i}{\|x_i\|}
\right\rangle
\right)\right)
\]

\textbf{Ключевое отличие} заключается в том, что в логистическом функционале отступ равен $M_i = y_i\langle w, x_i\rangle$ и зависит от нормы $w$, а в сферическом функционале отступ равен $M_i = \beta\, y_i \cos\theta(w, x_i)$ и \textbf{не зависит} от нормы вектора весов $\|w\|$ и нормы объектов $\|x_i\|$. Функционал \emph{инвариантен} относительно масштабирования $w$.
Роль «масштаба» отступа берёт на себя гиперпараметр $\beta$, сам отступ ограничен: $\cos\theta \in [-1,1]$, поэтому $M_i \in [-\beta, \beta]$\\[1mm]

\subsection*{2.}

Рассмотрим логистический функционал
\[
Q_{\mathrm{log}}(w) = \frac{1}{\ell}\sum_{i=1}^{\ell}
\log\bigl(1 + \exp(-y_i\langle w, x_i\rangle)\bigr)
\]
Для линейно разделимой выборки существует вектор $\hat{w}$ такой, что
$y_i\langle \hat{w}, x_i\rangle > 0$ для всех $i$. Тогда, рассматривая
$w = c\hat{w}$ при $c\to +\infty$, получаем:
\[
y_i\langle c\hat{w}, x_i\rangle = c \cdot y_i\langle \hat{w}, x_i\rangle
\to +\infty \quad \Longrightarrow \quad
Q_{\mathrm{log}}(c\hat{w}) \to 0
\]
Однако значение $Q_{\mathrm{log}} = 0$ \textbf{не достигается} ни при
каком конечном $w$ (поскольку $\log(1+e^{-z}) > 0$ при любом конечном $z$). Таким образом, инфимум функционала равен нулю, но минимум
\textbf{не достигается}: $\|w^*\| = +\infty$.\\[1mm]
При обучении градиентным спуском веса $w$ растут неограниченно
($\|w\|\to\infty$), что приводит к численным проблемам и отсутствию
единственного решения (направление $w/\|w\|$ стабилизируется, но сама
норма не ограничена)\\[1mm]
\textbf{В сферическом функционале} отступ равен $\beta\cos\theta(w,x_i)$
и не зависит от $\|w\|$\\[1mm] Масштабирование вектора весов
$w \mapsto cw$ ($c > 0$) не меняет значение функционала, так как
\[
\cos\theta(cw, x_i) = \frac{\langle cw, x_i\rangle}{\|cw\|\cdot\|x_i\|}
= \frac{c\langle w, x_i\rangle}{c\|w\|\cdot\|x_i\|}
= \cos\theta(w, x_i)
\]
Следовательно, увеличение нормы $w$ не уменьшает функционал.
Минимум ищется фактически по направлению $w/\|w\|$ на единичной сфере,
которая является компактным множеством. По теореме Вейерштрасса
непрерывная функция на компакте \textbf{достигает} своего минимума.

Таким образом, \textbf{проблема решается}: в сферическом функционале
минимум всегда достигается на конечном векторе весов, на
конечном направлении\\[1mm]

\subsection*{3.}

Параметр $\beta$ стоит как множитель перед $\cos\theta$:
\[
Q(w) = \frac{1}{\ell}\sum_{i=1}^{\ell}
\log\bigl(1 + \exp(-\beta\, y_i\cos\theta(w,x_i))\bigr)
\]
Поскольку $\cos\theta \in [-1, 1]$, то аргумент экспоненты лежит в
$[-\beta, \beta]$\\[1mm]
\begin{itemize}
    \item \textbf{$\beta \to 0$:}
    При малых $\beta$ аргумент экспоненты стремится к нулю для всех
    объектов:
    \[
    \log(1 + e^{-\beta y_i\cos\theta}) \approx
    \log(1 + 1 - \beta y_i\cos\theta) \approx \log 2
    \]
    Функционал становится практически константным ($Q \approx \log 2$)
    и перестаёт различать правильные и неправильные классификации.
    Градиенты становятся мал\'{ы}, оптимизация затрудняется.
    Функционал становится «нечувствительным» к выбору $w$.

    \item \textbf{$\beta \to +\infty$:}
    Потеря для правильно классифицированного объекта
    ($y_i\cos\theta > 0$):
    \[
    \log(1+e^{-\beta \cdot (\text{полож.})}) \to \log(1+0) = 0
    \]
    Потеря для неправильно классифицированного объекта
    ($y_i\cos\theta < 0$):
    \[
    \log(1+e^{\beta\cdot|\cdot|}) \approx \beta |y_i\cos\theta| \to +\infty
    \]
    Функционал начинает очень жёстко «штрафовать» неправильно
    классифицированные объекты и функция потерь приближается к
    хинжевой с масштабом $\beta$. Решение становится менее
    «мягким» --- классификатор стремится к полному разделению.
    При больших $\beta$ задача становится всё более похожей на
    максимизацию углового margin.
\end{itemize}

Таким образом, $\beta$ играет роль \textbf{температуры}: малые $\beta$
делают функционал «мягким» и нечувствительным, большие $\beta$ делают
его «жёстким» и чувствительным к ошибкам.

\section*{Задача 2.(2.5 балла)}
Ответьте на вопросы по решающим деревьям и композициям:

\begin{enumerate}
    \item Допустим, мы решили использовать следующую функцию хаотичности
    (impurity) при построении решающего дерева:
    \[
    H(R) = \max_{(x_i,y_i)\in R} y_i.
    \]
    Чем она плоха? Какие разбиения будут предпочитаться при обучении дерева?

    \item Представьте, что мы решили обучать градиентный бустинг так:
    \[
    \frac{1}{\ell}\sum_{i=1}^{\ell}(b_N(x_i) - s_i)^2 \to \min_{b_N};
    \qquad
    s_i = \sum_{n=1}^{N-1}
    \left.\frac{\partial L(y_i, z)}{\partial z}\right|_{z = b_n(x_i)}
    \]
    Чем плох такой подход? Попробуйте подумать, как будут меняться $s_i$
    по мере роста $N$. Исправьте приведённые выше формулы, чтобы
    получилась стандартная процедура градиентного бустинга.

    \item Известно, что градиентный бустинг является в некотором смысле
    аналогом градиентного спуска: каждая базовая модель приближает
    антиградиент ошибки на обучающей выборке. А как модифицировать
    бустинг, чтобы он стал аналогом mini-batch градиентного спуска?
    Запишите формулы для сдвигов и для обучения базовой модели в этом
    случае. Поясните, почему это правда будет аналог.
\end{enumerate}

\subsection*{1.}
Рассмотрим задачу регрессии или классификации с непрерывными метками.
Данная функция хаотичности равна максимальному значению целевой переменной
в узле $R$\\[1mm]
\textbf{Проблемы:}
\begin{enumerate}
    \item[] \textbf{Здесь нет учёта разнородности,} так как хорошая impurity-функция должна отражать разброс/неоднородность в узле (например, дисперсия для
    регрессии, энтропия/джини для классификации). Функция $\max y_i$
    зависит \emph{только от одного экстремального объекта} и не учитывает
    распределение целевых переменных остальных объектов

    \item[] При построении дерева
    мы минимизируем взвешенную сумму impurity дочерних узлов:
    \[
    \frac{|R_L|}{|R|}\,H(R_L) + \frac{|R_R|}{|R|}\,H(R_R) \to \min
    \]
    Подставляя $H(R) = \max y_i$:
    \[
    \frac{|R_L|}{|R|}\max_{(x_i,y_i)\in R_L} y_i +
    \frac{|R_R|}{|R|}\max_{(x_i,y_i)\in R_R} y_i \to \min
    \]
    Такая функция будет стремиться \textbf{отделить объекты с большими
    $y_i$ в как можно меньшее подмножество}. Объект с наибольшим $y_i$
    будет «вытеснен» в дочерний узел с наименьшим размером.
    Фактически дерево будет пытаться изолировать объекты-выбросы
    (с наибольшими $y$), а не строить осмысленное разбиение

    \item[] \textbf{Отсутствие симметрии,} весь функция не реагирует на
    минимальные значения $y$, а значит не штрафует за наличие
    разнородных (по $y$) объектов с малыми значениями

    \item[] \textbf{Не уменьшается при идеальном разбиении:} Если в узле
    значения $y$ одинаковы, $H = y_{\max}$ будет ненулевым, в отличие
    от дисперсии или энтропии, которые обнуляются при однородности.
    Это значит, что дерево не имеет стимула достигать однородности
    по $y$ в узлах
\end{enumerate}

\subsection*{2.}
Дана формулировка:
\[
\frac{1}{\ell}\sum_{i=1}^{\ell} (b_N(x_i) - s_i)^2 \to \min_{b_N};
\qquad
s_i = \sum_{n=1}^{N-1}
\left.\frac{\partial L(y_i, z)}{\partial z}\right|_{z = b_n(x_i)}
\]

В правильном градиентном бустинге на $N$-м шаге мы хотим, чтобы
очередной базовый алгоритм $b_N$ аппроксимировал \textbf{антиградиент}
функции потерь, вычисленный \textbf{в текущем приближении}
$a_{N-1}(x_i) = \sum_{n=1}^{N-1} b_n(x_i)$, а не сумму градиентов
в отдельных $b_n(x_i)$.

\textbf{Проблема 1:} Сдвиги $s_i$ вычисляются как \emph{сумма градиентов,
взятых в разных точках} $z = b_n(x_i)$, вместо одного градиента,
взятого в точке текущей композиции $z = a_{N-1}(x_i)$.

\textbf{Проблема 2:} Знак. Мы должны обучать базовую модель на
\emph{антиградиент} (со знаком минус).

\textbf{При росте $N$,} поскольку $s_i$ "--- сумма $N-1$ градиентов, её абсолютное значение
может неограниченно расти с ростом $N$. Базовый алгоритм $b_N$
(обычно дерево ограниченной глубины) не сможет аппроксимировать
такие растущие значения, и обучение будет деградировать.\\[1mm]
\textbf{Правильные формулы:}\\[1mm]
Пусть $a_{N-1}(x_i) = \sum_{n=1}^{N-1} b_n(x_i)$ "--- текущее
приближение, тогда сдвиги:
\[
s_i = -\left.\frac{\partial L(y_i, z)}{\partial z}
\right|_{z = a_{N-1}(x_i)}
= -\left.\frac{\partial L(y_i, z)}{\partial z}
\right|_{z = \sum_{n=1}^{N-1} b_n(x_i)}
\]
Базовая модель обучается приближать эти антиградиенты:
\[
b_N = \arg\min_{b}\; \frac{1}{\ell}\sum_{i=1}^{\ell}(b(x_i) - s_i)^2
\]
Обновление композиции с learning rate $\eta$:
\[
a_N(x) = a_{N-1}(x) + \eta \cdot b_N(x)
\]

\subsection*{3.}

В mini-batch градиентном спуске на каждом шаге градиент вычисляется не
по всей выборке, а по случайному подмножеству
$B_N \subset \{1,\ldots,\ell\}$, $|B_N| = m$\\[1mm]
\textbf{Аналог в бустинге --- стохастический градиентный бустинг:}\\[1mm]
На каждой итерации $N$
\begin{enumerate}
    \item Выбираем случайное подмножество
    $B_N \subset \{1,\ldots,\ell\}$, $|B_N| = m$

    \item Вычисляем сдвиги \textbf{только для объектов из $B_N$}:
    \[
    s_i = -\left.\frac{\partial L(y_i, z)}{\partial z}
    \right|_{z = a_{N-1}(x_i)}, \quad i \in B_N
    \]

    \item Обучаем базовую модель \textbf{только по объектам из $B_N$}:
    \[
    b_N = \arg\min_{b}\; \frac{1}{|B_N|}\sum_{i \in B_N}(b(x_i) - s_i)^2
    \]

    \item Обновляем:
    \[
    a_N(x) = a_{N-1}(x) + \eta\, b_N(x)
    \]
\end{enumerate}

В обычном градиентном бустинге мы аппроксимируем антиградиент
функционала $Q = \frac{1}{\ell}\sum_{i=1}^{\ell} L(y_i, a(x_i))$
по всем объектам --- это аналог полного (batch) градиентного спуска
в пространстве значений модели на обучающей выборке.

Если же мы берём случайное подмножество $B_N$, то мы аппроксимируем
\emph{стохастический антиградиент}
\[
-\frac{1}{|B_N|}\sum_{i\in B_N}
\left.\frac{\partial L(y_i,z)}{\partial z}\right|_{z=a_{N-1}(x_i)}
\]
который является несмещённой оценкой полного антиградиента\\[1mm]
Это в точности \textbf{аналогия с mini-batch SGD}, где вместо полного градиента
используется его оценка по подвыборке. Это вносит стохастичность,
которая действует как регуляризация и ускоряет обучение.

\section*{Задача 3. (2.5 балла)}

Тимофей поручил студентам специализации «Машинное обучение и приложения»
разработать модель для классификации состояния стойл на ферме на чистые и
нуждающиеся в уборке. Стоимость работ по уборке стойла составляет 200 рублей.
Если стойло ошибочно было классифицировано как не нуждающееся в уборке,
ферма теряет скот и несёт убытки в 1000 рублей.

В обучающей выборке дано 8 стойл, для которых моделью были выданы следующие
вероятности того, что они нуждаются в ремонте или уборке:

\begin{center}
\begin{tabular}{c|cccccccc}
$b(x)$ & 0.1 & 0.1 & 0.4 & 0.4 & 0.4 & 0.8 & 0.8 & 0.95 \\
\midrule
$y$     & 0   & 0   & 0   & 1   & 0   & 1   & 0   & 1
\end{tabular}
\end{center}

Выполните задания:

\begin{enumerate}
    \item По обучающей выборке подберите порог классификации, который
    соответствует наименьшим издержкам фермы. Рассчитайте Accuracy и
    F-меру для обучающей выборки, исходя из выбранного вами порога.

    \item Вычислите AUC-ROC для $b(x)$ на обучающей выборке. Если у
    нескольких объектов выход модели совпадает и равен $p$, то в
    ROC-кривой соединяйте отрезком точки, соответствующие порогам
    $p - \varepsilon$ и $p + \varepsilon$.

    \item Допустим теперь, что мы имеем $K$ состояний стойла. Запишите
    функцию потерь для этой задачи многоклассовой классификации, и
    предложите, как сделать так, чтобы в ней учитывались различные
    издержки $c_{km}$ за выдачу класса $m$ при правильном классе $k$.

    \item Предложите обобщение AUC-ROC для задачи классификации с тремя
    классами. Объясните, как вычислять такую метрику.

    \emph{Подсказка:} аналитическая формула AUC-ROC выглядит так:
    \[
    \frac{1}{\ell_-\ell_+}\sum_{\substack{i,j:\\y_i=1,\,y_j=0}}
    \left(\ind[b(x_i) > b(x_j)] +
    \tfrac{1}{2}\ind[b(x_i) = b(x_j)]\right)
    \]
    где $\ell_-, \ell_+$ "--- количество объектов отрицательного и
    положительного классов. Попробуйте придумать, как в этой формуле
    изменить множество, по которому идёт суммирование, и как поменять
    нормировку, чтобы формула подходила для многоклассового случая.
\end{enumerate}

\subsection*{1.}

При пороге $t$ мы предсказываем $\hat{y}_i = \ind[b(x_i) \ge t]$

\begin{itemize}
    \item Положительный класс ($\hat{y}=1$): стойло нуждается в уборке
    \item FP: стойло чистое, но мы убираем $\Rightarrow$
    издержки 200 руб
    \item FN: стойло грязное, но мы не убираем $\Rightarrow$
    убытки 1000 руб
    \item TP и TN --- без дополнительных ошибочных издержек, ведь TP правильная уборка --- можно считать что стоимость уборки
    неизбежна, мы минимизируем \emph{ошибочные} издержки
\end{itemize}
Ожидаемые ошибочные издержки:
\[
C(t) = 200 \cdot \text{FP}(t) + 1000 \cdot \text{FN}(t)
\]
Рассмотрим все осмысленные пороги:\\[1mm]
\textbf{Порог $t > 0.95$} (все предсказания 0):\\[1mm]
TP=0, FP=0, FN=3 (пропускаем все 3 грязных), TN=5 $\qquad C = 200\cdot 0 + 1000\cdot 3 = 3000$\\[1mm]
\textbf{Порог $t = 0.95$} ($b(x)\ge 0.95 \Rightarrow$ предсказываем 1):\\[1mm]
Объекты с $b(x)\ge 0.95$: объект 8 ($y=1$): $\hat{y}=(0,0,0,0,0,0,0,1) \\[1mm]
$ TP=1, FP=0, FN=2, TN=5, $\quad C = 200\cdot 0 + 1000\cdot 2 = 2000$\\[1mm]
\textbf{Порог $t = 0.8$} ($b(x)\ge 0.8$):\\[1mm]
Объекты: 6($y=1$), 7($y=0$), 8($y=1$). $\hat{y}=(0,0,0,0,0,1,1,1)$\\[1mm]
TP=2, FP=1, FN=1, TN=4, $\quad C = 200\cdot 1 + 1000\cdot 1 = 1200$\\[1mm]
\textbf{Порог $t = 0.4$} ($b(x)\ge 0.4$):\\[1mm]
Объекты: 3($y=0$), 4($y=1$), 5($y=0$), 6($y=1$), 7($y=0$), 8($y=1$)\\[1mm]
TP=3, FP=3, FN=0, TN=2, $\quad C = 200\cdot 3 + 1000\cdot 0 = 600$\\[1mm]
\textbf{Порог $t = 0.1$} ($b(x)\ge 0.1$):\\[1mm]
Все объекты предсказаны как 1:\\[1mm]
TP=3, FP=5, FN=0, TN=0, $ C = 200\cdot 5 + 1000\cdot 0 = 1000$\\[1mm]
\textbf{Вывод:} минимальные издержки $C = 600$ достигаются при пороге
$t = 0.4$\\[1mm]
\textbf{ Accuracy и F-мера при $t = 0.4$:}\\[1mm]
Confusion matrix: TP=3, FP=3, FN=0, TN=2
\[
\text{Accuracy} = \frac{TP+TN}{TP+TN+FP+FN} = \frac{3+2}{3+2+3+0}
= \frac{5}{8} = 0.625
\]
\[
\text{Precision} = \frac{TP}{TP+FP} = \frac{3}{3+3} = \frac{1}{2} = 0.5
\]
\[
\text{Recall} = \frac{TP}{TP+FN} = \frac{3}{3+0} = 1
\]
\[
\text{F}_1 = 2\cdot\frac{\text{Precision}\cdot\text{Recall}}
{\text{Precision}+\text{Recall}}
= 2\cdot\frac{0.5\cdot 1}{0.5 + 1} = \frac{1}{1.5}
= \frac{2}{3} \approx 0.667
\]

\subsection*{2.}

Для построения ROC-кривой перебираем пороги от $+\infty$ до $-\infty$ и
для каждого вычисляем FPR и TPR. В выборке: $\ell_+ = 3$, $\ell_- = 5$\\[1mm]
Уникальные значения $b(x)$: $0.95, 0.8, 0.4, 0.1$\\[1mm]
По условию, при совпадающих значениях $b(x)$ мы соединяем точки для
порогов $p-\varepsilon$ и $p+\varepsilon$ отрезком.

\begin{tikzpicture}[scale=7]

\draw[->] (-0.05,0) -- (1.15,0) node[right] {FPR};
\draw[->] (0,-0.05) -- (0,1.15) node[above] {TPR};

\draw[gray, dashed, thin] (0,0) -- (1,1);
\node[gray, font=\scriptsize, rotate=45] at (0.72,0.65) {случайный};

\draw[blue!70!black, very thick]
    (0, 0) -- (0, 0.3333)
            -- (0.2, 0.6667)
            -- (0.6, 1.0)
            -- (1.0, 1.0);

\filldraw[red!80!black] (0, 0)       circle (0.8pt);
\filldraw[red!80!black] (0, 0.3333)  circle (0.8pt);
\filldraw[red!80!black] (0.2, 0.6667) circle (0.8pt);
\filldraw[red!80!black] (0.6, 1.0)   circle (0.8pt);
\filldraw[red!80!black] (1.0, 1.0)   circle (0.8pt);

\node[right, font=\scriptsize, text=red!80!black] at (0.02, 0.0)
    {$(0,\;0)$\; $t>0.95$};
\node[right, font=\scriptsize, text=red!80!black] at (0.02, 0.3333)
    {$(0,\;\frac{1}{3})$\; $t=0.95$};
\node[above right, font=\scriptsize, text=red!80!black] at (0.2, 0.6667)
    {$(\frac{1}{5},\;\frac{2}{3})$\; $t=0.8$};
\node[above right, font=\scriptsize, text=red!80!black] at (0.6, 1.0)
    {$(\frac{3}{5},\;1)$\; $t=0.4$};
\node[below left, font=\scriptsize, text=red!80!black] at (1.0, 1.0)
    {$(1,\;1)$\; $t=0.1$};

\fill[blue!10, opacity=0.5]
    (0,0) -- (0, 0.3333) -- (0.2, 0.6667) -- (0.6, 1.0) -- (1.0, 1.0)
    -- (1.0, 0) -- cycle;

\node[blue!50!black, font=\footnotesize] at (0.55, 0.35) {AUC};
\end{tikzpicture}

\textbf{AUC методом трапеций:}

\[
\text{AUC-ROC} = 0 + \frac{1}{10} + \frac{1}{3} + \frac{2}{5}
= \frac{3 + 10 + 12}{30} = \frac{25}{30} = \frac{5}{6} \approx 0.833
\]

\textbf{Проверка аналитической формулой:}
\[
\text{AUC} = \frac{1}{\ell_-\ell_+}\sum_{\substack{i,j:\\y_i=1,\,y_j=0}}
\left(\ind[b(x_i)>b(x_j)] + \tfrac{1}{2}\ind[b(x_i)=b(x_j)]\right)
\]

Положительные объекты: 4($b=0.4$), 6($b=0.8$), 8($b=0.95$)\\[1mm]
Отрицательные объекты: 1($b=0.1$), 2($b=0.1$), 3($b=0.4$),
5($b=0.4$), 7($b=0.8$)\\[1mm]

$\ell_+\cdot\ell_- = 3\cdot 5 = 15$ пар. Считаем вклад каждой:

\begin{center}
\begin{tabular}{c|ccccc|c}
$(+)\backslash(-)$ & 1(0.1) & 2(0.1) & 3(0.4) & 5(0.4) & 7(0.8) & сумма\\
\midrule
4(0.4) & 1 & 1 & 1/2 & 1/2 & 0 & 3\\
6(0.8) & 1 & 1 & 1 & 1 & 1/2 & 4.5\\
8(0.95) & 1 & 1 & 1 & 1 & 1 & 5\\
\midrule
& & & & & & 12.5
\end{tabular}
\end{center}

\[
\text{AUC} = \frac{12.5}{15} = \frac{25}{30}
= \frac{5}{6} \approx 0.833
\]

\subsection*{3.}

Пусть имеется $K$ классов. Модель выдаёт вероятности
$p_k(x) = P(y=k|x)$ для $k=1,\ldots,K$.
Введём матрицу издержек $C = (c_{km})_{K\times K}$, где $c_{km}$ "---
стоимость ошибки при предсказании класса $m$, когда истинный класс $k$.

\textbf{Стандартная кросс-энтропия:}
\[
L(y, p) = -\sum_{k=1}^{K} \ind[y=k]\log p_k(x)
= -\log p_y(x)
\]
Из-за издержек нужна модификация, например, есть подход --- использовать \textbf{cost-sensitive}
кросс-энтропию с весами, зависящими от класса и ошибки:
\[
L(y, p) = -\sum_{k=1}^{K} \ind[y=k] \sum_{m=1}^{K} c_{km}\, p_m(x)
\]
Однако более стандартный подход --- учитывать издержки на этапе
принятия решения. Модель обучается предсказывать вероятности
$p_k(x)$, а \emph{решающее правило} строится так:
\[
\hat{y}(x) = \arg\min_{m=1,\ldots,K}
\sum_{k=1}^{K} c_{km}\, p_k(x) = \arg\min_m E_{y|x}[c_{y,m}]
\]
То есть мы предсказываем тот класс, который минимизирует ожидаемые
издержки.

Альтернативный подход --- \textbf{взвешенная кросс-энтропия}:
\[
L(y, p) = -w_y \log p_y(x)
\]
где $w_k = \sum_{m \ne k} c_{km}$ "--- суммарная стоимость ошибок
для класса $k$\\[1mm]

\subsection*{4.}

Для бинарной классификации AUC-ROC --- это вероятность того, что случайный
положительный объект будет ранжирован выше случайного отрицательного:
\[
\text{AUC} = \frac{1}{\ell_-\ell_+}\sum_{\substack{i,j:\\y_i=1,y_j=0}}
\left(\ind[b(x_i)>b(x_j)] + \tfrac{1}{2}\ind[b(x_i)=b(x_j)]\right)
\]

\textbf{Обобщение для $K=3$ классов}\\[1mm]
Используем подход \textbf{one-vs-one}: для каждой пары классов
$(k, m)$, $k \ne m$, вычисляем AUC, рассматривая задачу «отличить
класс $k$ от класса $m$» с использованием скоринга $b_k(x) - b_m(x)$
(или $b_k(x)$, если модель выдаёт скоры для каждого класса)
\[
\text{AUC}_{k,m} = \frac{1}{\ell_k\,\ell_m}
\sum_{\substack{i,j:\\y_i=k,\,y_j=m}}
\left(\ind [b_k(x_i) > b_k(x_j)] +
\tfrac{1}{2}\ind [b_k(x_i) = b_k(x_j)]\right)
\]
Итоговая метрика --- усреднение по всем парам классов:
\[
\text{AUC}_{\text{multi}} = \frac{2}{K(K-1)}
\sum_{k<m} \text{AUC}_{k,m} =
\frac{1}{3}\bigl(\text{AUC}_{1,2} + \text{AUC}_{1,3} + \text{AUC}_{2,3}\bigr)
\]
\textbf{Как изменить аналитическую формулу:}
\begin{itemize}
    \item Множество суммирования: вместо пар $(i,j)$ с $y_i=1, y_j=0$
    берём все пары $(i,j)$ с $y_i \ne y_j$ (или фиксируем пару классов
    и суммируем по объектам этих двух классов).
    \item Нормировка: вместо $\ell_-\ell_+$ "--- для каждой пары классов
    $(k,m)$ нормируем на $\ell_k\cdot\ell_m$, затем усредняем по
    $\binom{K}{2}$ парам.
\end{itemize}

Альтернатива "--- подход \textbf{one-vs-all}: для каждого класса $k$
вычислять AUC как «класс $k$ vs остальные» и усреднять.


\section*{Задача 4. (2.5 балла)}

Пусть логарифм обменного курса в день $t$ изменяется на валютном рынке
согласно закону
\[
r_t = 0.8\,r_{t-1} + 0.6\,\varepsilon_t \qquad r_0 = \varepsilon_0
\]
Все случайные величины $\varepsilon_t$ независимы и имеют распределение
$\mathcal{N}(0,1)$

Биржевой аналитик Одэхингум знает этот закон и для построения модели
предсказания курса на следующий день использует метод обучения
\[
\mu(X)(r_{t-1}) = 0.8\,r_{t-1} + 0.6\,\gamma,
\]
где число $\gamma$ семплируется из стандартного нормального распределения
независимо при каждом обучении модели. Найдите шум, смещение и разброс
метода обучения $\mu$ при предсказании $r_2$.

Для вашего удобства напомним формулу разложения ошибки:
\[
L(\mu) = \underbrace{E_{x,y}\bigl[(y - E[y\mid x])^2\bigr]}_{\text{шум}}
+ \underbrace{E_x\bigl[(E_X[\mu(X)] - E[y\mid x])^2\bigr]}_{\text{смещение}}
+ \underbrace{E_x\bigl[E_X[(\mu(X) - E_X[\mu(X)])^2]\bigr]}_{\text{разброс}}
\]

\textbf{Раскроем}
\begin{align}
r_0 &= \varepsilon_0, \\
r_1 &= 0.8\,\varepsilon_0 + 0.6\,\varepsilon_1, \\
r_2 &= 0.8\,r_1 + 0.6\,\varepsilon_2
= 0.8(0.8\varepsilon_0 + 0.6\varepsilon_1) + 0.6\varepsilon_2
= 0.64\,\varepsilon_0 + 0.48\,\varepsilon_1 + 0.6\,\varepsilon_2.
\end{align}

Здесь в контексте Bias-Variance разложения: $x = r_1$ "--- входной признак, $y = r_2$ "--- целевая переменная, $X$ "--- обучающая выборка, здесь определяется значением $\gamma$\\[1mm]

\textbf{Условное математическое ожидание $E[y|x] = E[r_2|r_1]$}

\[
r_2 = 0.8\,r_1 + 0.6\,\varepsilon_2
\]
Поскольку $\varepsilon_2$ независим от $r_1$:
\[
E[r_2 | r_1] = 0.8\,r_1 + 0.6\cdot E[\varepsilon_2]
= 0.8\,r_1
\]

\textbf{Шум}

\[
\text{Noise} = E_{x,y}\bigl[(y - E[y|x])^2\bigr]
= E\bigl[(r_2 - 0.8\,r_1)^2\bigr]
= E\bigl[(0.6\,\varepsilon_2)^2\bigr]
= 0.36\,E[\varepsilon_2^2] = 0.36
\]

\textbf{Ожидание предсказания $E_X[\mu(X)(r_1)]$}

Обучающая выборка «определяется» значением $\gamma$, поэтому
$E_X[\cdot] = E_\gamma[\cdot]$:
\[
E_\gamma[\mu(X)(r_1)] = E_\gamma[0.8\,r_1 + 0.6\,\gamma]
= 0.8\,r_1 + 0.6\cdot E[\gamma] = 0.8\,r_1
\]

\textbf{Смещение}

\[
\text{Bias} = E_x\bigl[(E_X[\mu(X)(x)] - E[y|x])^2\bigr]
= E_{r_1}\bigl[(0.8\,r_1 - 0.8\,r_1)^2\bigr]
= E_{r_1}[0] = 0
\]

\textbf{ Разброс}

\[
\text{Variance} = E_x\bigl[E_X\bigl[(\mu(X)(x) - E_X[\mu(X)(x)])^2
\bigr]\bigr]
\]

Вычислим внутреннее ожидание:
\[
\mu(X)(r_1) - E_\gamma[\mu(X)(r_1)]
= (0.8\,r_1 + 0.6\gamma) - 0.8\,r_1 = 0.6\gamma
\]

\[
E_\gamma\bigl[(0.6\gamma)^2\bigr] = 0.36\,E[\gamma^2] = 0.36
\]

Это не зависит от $r_1$, поэтому:
\[
\text{Variance} = E_{r_1}[0.36] = 0.36
\]

\textbf{Итого:}

\[
L(\mu) = \underbrace{0.36}_{\text{шум}} +
\underbrace{0}_{\text{смещение}} +
\underbrace{0.36}_{\text{разброс}} = 0.72
\]

\begin{itemize}
    \item \textbf{Шум} = 0.36 "--- это неустранимая ошибка, обусловленная
    стохастичностью $\varepsilon_2$ в модели генерации данных.
    \item \textbf{Смещение} = 0 "--- в среднем по обучающим выборкам (по
    $\gamma$) предсказание метода совпадает с $E[r_2|r_1]$, так как
    $E[\gamma] = 0$. Метод является несмещённым.
    \item \textbf{Разброс} = 0.36 "--- при разных обучениях (разных $\gamma$)
    предсказание отклоняется на $0.6\gamma$ от среднего. Дисперсия этого
    отклонения равна шуму, что логично: метод «подменяет» истинный шум
    $\varepsilon_2$ независимой копией $\gamma$ с тем же распределением.
\end{itemize}

\end{document}